% Part: incompleteness
% Chapter: incompleteness-provability
% Section: introduction

\documentclass[../../../include/open-logic-section]{subfiles}

\begin{document}

\olfileid{inc}{inp}{int}

\olsection{پېژندنه}

هېلبرټ په دې اند و چې د يو رياضيکي جوړښت، لکه د طبيعي عددونو، دپاره
د بديهي اصولو نظام تر هغه وخته بسيا نۀ دے چې د هماغه جوړښت ټولې
رښتينې جملې ترې استنتاج نۀ شي۔ د صوري استنتاج له نظامونو سره د هغه
وروستۍ لېوالتيا چې ورسره يوځاے کړو، نو داسې ښکاري چې هغه غوښتل د
طبيعي عددونو په باب زموږ صوري استدلالي نظامونه يوازې سازګار نۀ، بلکې
\emph{بشپړ} هم وي؛ يعنې د نظام د ژبې هره جمله يا پخپله !!{derivable}
وي يا ئې نفي همداسې وي۔ د ګوډل د ناتکميلۍ لومړۍ قضيه ښيي چې
داسې د بديهي اصولو نظام نشته: د حساب دپاره هېڅ بشپړ، سازګار او
!!{axiomatizable} صوري نظام نشته۔ په حقيقت کښې، هېڅ «په کافي اندازه
پياوړې»، سازګاره او !!{axiomatizable} رياضيکي تيوري بشپړه نۀ ده۔

د هېلبرټ يو لا مهم هدف، چې د نوې («کلاسيکې») رياضي د توجيه دپاره
د هغه د پروګرام منځنے ټکی و، د کلاسيک استدلال استازيتوب کوونکو صوري
نظامونو د سازګارۍ متناهي‌پاله ثبوتونه موندل وو۔ نو د هېلبرټ د
پروګرام له پلوه د ګوډل د ناتکميلۍ دويمه قضيه تر لومړۍ ډېر دروند ګوزار
و۔ دويمه قضيه هم د لومړۍ په شان په لږ ناڅرګندو ټکو بيانېدای شي:
نږدې معنا ئې دا ده چې د حساب هېڅ په کافي اندازه پياوړې تيوري خپله
سازګاري نۀ شي ثابتولے۔ د «په کافي اندازه پياوړې» مطلب به له~$\Th{Q}$
څخه لږ څه زيات وغواړي۔

د ناتکميلۍ د قضيې د ګوډل د اصلي ثبوت مفکوره د اپيمنيدس په تناقض
کښې موندل کېږي۔ اپيمنيدس، چې د کريت اوسېدونکے و، وويل چې د کريت
ټول اوسېدونکي دروغجن دي؛ د دې تناقض لا نېغه بڼه «دا جمله ناسمه ده»
ويل دي۔ په اصل کښې ګوډل د رښتياوالي پر ځاے !!{derivability} وکاروله
او !!a{sentence} ئې صوري کړه چې په غير مستقيم ډول وايي: زه پخپله
!!{derivable} نۀ يم۔ که دغه !!{sentence} !!{derivable} واے، تيوري به
ناسازګاره وه۔ ګوډل دا هم وښودله چې د هغه د بديهي اصولو له نظامه د
دغې !!{sentence} نفي هم !!{derivable} نۀ ده۔ (د دې دويمې برخې دپاره
ګوډل اړ و چې فرض کړي تيوري~$\Th{T}$ «$\omega$-سازګاره» ده۔
$\omega$-سازګاري له سازګارۍ سره تړاو لري، خو تر هغې پياوړے خاصيت
دے۔\footnote{يعنې هره $\omega$-سازګاره تيوري سازګاره ده، خو
برعکس لازمه نۀ ده۔} له ګوډل څو کاله وروسته
روسر وښودله چې يوازې د~$\Th{T}$ عادي سازګاري فرض کول بس دي۔)

لومړۍ ستونزه دا پوهېدل دي چې داسې !!a{sentence} څنګه جوړېږي چې
ځان ته مراجعه وکړي۔ د~$\Th{Q}$ په ژبه کښې د هر !!{formula}~$!A$
دپاره دې~$\gn{!A}$ هغه عدديه نښه وي چې له~$\Gn{!A}$ سره سمون خوري۔
د دې معنا ته ځير شئ: $!A$~د~$\Th{Q}$ په ژبه کښې !!a{formula} دے؛
$\Gn{!A}$~يو طبيعي عدد دے؛ او $\gn{!A}$~د~$\Th{Q}$ په ژبه کښې يو
\emph{ترم} دے۔ نو د~$\Th{Q}$ د ژبې هر !!{formula}~$!A$ يو
\emph{نوم}، يعنې~$\gn{!A}$، لري چې د~$\Th{Q}$ په ژبه کښې ترم
دے۔ دا موږ ته
هغه مفهومي چوکاټ راکوي چې په هغه کښې د~$\Th{Q}$ د ژبې
!!{formula}s د نورو !!{formula}s په باب څه «ويلے» شي۔ لاندې لمه
د ثابت ټکي د لمې په نوم يادېږي۔

\begin{lem}
فرض کړئ $\Th{T}$ هره هغه تيوري ده چې~$\Th{Q}$ غځوي، او $!B(x)$
هر هغه !!{formula} دے چې يوازې متغير~$x$ پکې آزاد وي۔ نو داسې
!!a{sentence}~$!A$ شته چې
$\Th{T} \Proves!A \liff !B(\gn{!A})$۔
\end{lem}

لمه وايي چې د هر خاصيت~$!B(x)$ دپاره داسې !!a{sentence}~$!A$
شته چې وايي «$!B(x)$ زما په باب رښتيا دے»، او~$\Th{T}$ هم دغه
برابري «پېژني»۔

داسې !!a{sentence} څنګه جوړولے شو؟ د اپيمنيدس د تناقض لاندې بڼه
وګورئ چې کواین وړاندې کړې ده:
\begin{quote}
«که يې نقل‌قول ترې وړاندې کېښودل شي، دروغ وايي» که يې نقل‌قول
ترې وړاندې کېښودل شي، دروغ وايي۔
\end{quote}
دا جمله په نېغه ځان ته مراجعه نۀ کوي۔ يوازې د نقل‌قول د نښو ترمنځ
د نحوي شيانو په باب خبره کوي، نو له دې جملو سره په يوه کچه ده:
\begin{enumerate}
\item «رابرټ» يو ښۀ نوم دے۔
\item «زه وځغاستم۔» يوه لنډه جمله ده۔
\item «درې کلمې لري» درې کلمې لري۔
\end{enumerate}
خو که «که يې نقل‌قول ترې وړاندې کېښودل شي، دروغ وايي» عبارت واخلو
او د همدغه عبارت نقل شوې بڼه ترې مخکې کېږدو، څه پېښېږي؟ هماغه
لومړۍ جمله ترې جوړېږي!{} په لنډه توګه، جمله وايي چې پخپله ناسمه ده۔

\end{document}
